\documentclass[11pt]{article}
\usepackage[margin=1.15in]{geometry}
\usepackage[T1]{fontenc}
\usepackage{lmodern}
\usepackage{microtype}
\usepackage{amsmath,amssymb,amsthm}
\usepackage{setspace}
\usepackage[hidelinks]{hyperref}
\onehalfspacing

\newtheorem{proposition}{Proposition}
\theoremstyle{definition}
\newtheorem{definition}{Definition}

\title{Sufficiency Without Standing: Corporate Law and the Sale of a Public Square}
\author{Flyxion\\Independent Researcher}
\date{November 2026}

\begin{document}
\maketitle

\begin{abstract}
\noindent Ann Lipton's \emph{Every Billionaire Is a Policy Failure} reads the 2022 acquisition of Twitter as a case in which corporate law performed exactly as designed and still produced a public failure. This essay restates that result in a constraint-first vocabulary. A sale is treated as a transition between states of a system, and a legal procedure is treated as a projection that retains some features of the transition and discards others. The projection used by Delaware merger law is sufficient for the question it was built to answer, which is whether the selling shareholders received the agreed price, and it is silent on the question the public needed answered, which is who is authorized to hold the power that the sale transfers. Shareholder primacy is then analysed as a division of labour in which the sale procedure gave shareholders a recognized role while leaving the people subject to the transferred powers without a comparable avenue to contest the terms of transfer. The wealth that financed the transition is analysed as the output of an earlier process in which enforcement lagged, and a possible feedback between concentrated wealth and attenuated oversight is described as a loop rather than an event. Authority to complete a sale and legitimacy to exercise the powers it transfers are distinct, and a legitimate transfer of a communications platform would have to assign standing and the obligations of stewardship at the point where the capacities change hands.
\end{abstract}

\section{Introduction}

The usual account of the Twitter acquisition is a story about a person. A wealthy buyer made an offer, regretted it, attempted to withdraw, was sued, and eventually closed at the agreed price of \$54.20 per share. Lipton's article declines to tell it that way. Her subject is the institutional arrangement in which the story could occur at all, and her finding is uncomfortable because it does not depend on anyone behaving badly \cite{lipton2024}. The board accepted an offer of \$54.20 per share, a price that became a windfall for shareholders, an overpayment relative to standalone value by as much as \$24 billion on the estimates Lipton cites, once the market fell. When the buyer tried to leave, Twitter sued to hold him to the bargain, the rulings of the Court of Chancery made clear that he held the weaker hand, and he abandoned his opposition before the court decided what remedy to order. By the standards that corporate law applies to its own performance, every step was a success. By any standard that attends to the people who used the platform, worked at it, reported through it, or built communities on it, the same sequence delivered an important communications system into the hands of a single owner without any procedure in which those people could register a claim.

The argument developed here accepts that finding and asks what kind of object a legal procedure must be for success and failure to coexist in this way. The answer proposed is structural. A procedure is a projection of a larger transition onto a smaller set of variables, and a projection can be sufficient for one set of questions while being blind to another. Corporate law projected the sale of Twitter onto the variables of price, premium, and enforceability of the agreement, and on those variables the outcome was correct. The variables on which the public's stake depended, namely the governance of speech, the conditions of employment, and the continued availability of a shared infrastructure for communication, were not represented in the projection at all. The failure therefore does not live inside the procedure. It lives in the gap between what the procedure retains and what the transition changes.

This framing has a consequence for how reform should be imagined. If the problem were that the court reached a wrong result, the remedy would be a better rule for the court. If the problem is that the projection omits the relevant variables, then no refinement of the rule within the projection can help, because the rule never sees what it would need to weigh. Lipton reaches a version of this conclusion when she doubts that allowing Delaware to weigh harms to outsiders would address the underlying concentration of power. The vocabulary of projections and standing explains why that doubt is well founded.

\section{A Sale as a Transition}

A constraint-first approach begins from what a system is permitted to do and treats the changes it undergoes as primary, with states and parties appearing as the endpoints of transitions. On that approach the sale of Twitter is a single transition with a before and an after. Before the transition, the platform was governed by a board answerable to dispersed shareholders whose interests were largely financial and whose individual influence over product decisions was negligible. After it, the platform was governed by one owner whose discretion was constrained chiefly by the law that applies to any business and by the commercial need to retain users and advertisers. Between the two states, a set of capacities changed hands: the power to decide what is amplified, the power to admit or exclude accounts, the power to set conditions of work for the employees who maintain the system, and the power to alter the interfaces on which third parties had built.

Describing the transaction this way makes visible something that the price-centred description conceals. The transition is large, and most of its content is not about money. The consideration paid to shareholders is a transfer between two parties that is complete and measurable, but it is the smallest part of what changed, because the capacities that changed hands are exercised on and over a population that is much larger than either the buyer or the sellers. The transition has many affected parties and only two who were party to it.

\begin{definition}
A \emph{standing relation} holds between a party and a transition when the party is entitled to have its stake in the transition represented in the procedure that authorizes it, relative to a stated consequence, a stated scope, and a stated procedure of evaluation.
\end{definition}

The definition is deliberately relational. Standing is not a property that a person carries with them into every context. It is assigned relative to a consequence, which here is the transfer of control over a platform, relative to a scope, which here is the population the platform reaches, and relative to an evaluation procedure, which here is whatever mechanism approves or enforces the sale. Under this definition the selling shareholders had standing and the users did not, and the interesting question is whether that assignment tracks the stake each party had in the consequence. For the selling shareholders, the stake as shareholders was concluded at closing, since cash replaced their holding and nothing that happened afterward touched them as holders, even if some of them remained users or employees. Users and employees already depended on the platform before the sale. What began at closing was their exposure to the new controller's exercise of the powers that changed hands, and that exposure has no natural end.

\section{The Shareholder Projection and Its Sufficiency}

Shareholder primacy can be stated as a claim about projections. Managers are directed to evaluate transitions on one variable, the wealth of the shareholders, and the law supplies incentives that keep their attention there. The doctrine is old and has been argued in both directions, from the statement in \emph{Dodge v.\ Ford Motor Co.}\ that a business corporation is organized and carried on primarily for the profit of its stockholders \cite{dodge1919}, to the later and more polemical claim that the social responsibility of business is to increase its profits \cite{friedman1970}. What matters for the present argument is not whether the doctrine is wise but what makes it defensible on its own terms. Its defence is a division of labour. Corporate law is allowed to be single-minded because other bodies of law are assigned the protection of everyone else: labour law protects workers, safety regulation protects consumers and the public, securities regulation protects investors from deception, and competition law protects markets. The projection is justified because the complement of the projection is covered elsewhere. The division was always a rough approximation and always contested, as the long debate over the purposes of the corporation shows, from Berle's argument that corporate powers are held in trust for shareholders \cite{berle1931} and Dodd's reply that managers are trustees for a wider set of interests \cite{dodd1932} onward. Twitter is therefore not the first test of the approximation, and the claim here is that it is a case in which the approximation fails visibly and at scale.

Call this the sufficiency claim. It asserts that the shareholder projection is sufficient for the purposes of corporate law because the variables it discards are retained by other projections that together cover the whole. Sufficiency is always sufficiency for a class of questions and under a class of future revisions. A projection that suffices when the transitions it evaluates are small relative to the surrounding regulatory structure may fail when a transition is large relative to that structure, because the regulations that were meant to cover the discarded variables were written for an environment in which no single owner could alter the conditions under which they operate.

Omission alone does not establish a defect. A variable left out of a procedure may be irrelevant to the decision, may be fixed by a background constraint, or may be inferable from the variables the procedure retains, and in each of these cases the omission costs nothing. The argument therefore needs a criterion for when omission matters. Let $\pi(T)$ denote the information that a procedure retains about a possible transition $T$, and let $L(T)$ denote the status of $T$ under a specified public authorization criterion, for example whether the transfer preserves access, continuity, and accountability to a stated standard.

\begin{proposition}
Let $T_1$ and $T_2$ be admissible transitions with $\pi(T_1)=\pi(T_2)$ and $L(T_1)\neq L(T_2)$. Then no rule that operates on $\pi$ alone can determine $L$ correctly in both cases. A procedure is sufficient for the public criterion only if $L$ factors through $\pi$, that is, only if $L(T)$ is a function of $\pi(T)$ over the class of admissible transitions.
\end{proposition}

The proposition follows from its hypotheses once they are granted, and it is stated for the sake of the condition it isolates. The operative question is whether $L$ factors through $\pi$ over the class of transitions the procedure actually meets, and that is an empirical and institutional matter that the proposition does not settle. It is an informational limit and nothing more. Two sales can be identical in price, premium, and enforceability of the agreement while differing in whether the buyer is bound to preserve access or maintain continuity, and a rule that sees only price and enforceability must treat them alike. The proposition does not say that the court erred, and it does not say that affected parties are entitled to a voice. The first claim concerns what a procedure can determine. The second is a normative claim about who may contest what the procedure determines, and it requires its own defence, which is supplied by the conditions of scope matching, remainder assignment, and revisability developed in Section 7. The two claims come apart in a way that matters for the case. A procedure can be made to see a consequence, for instance by admitting evidence about the effects of a sale on users and employees, without those groups acquiring any entitlement to participate, to seek conditions, or to request later review. Visibility is a property of the procedure, and standing is a relation between the procedure and a party. The reforms that address the first do not automatically supply the second.

For a transfer of control over a communications platform, the complementary procedures that might have covered the omitted variables are thin. Labour law protects employees against specific wrongs, such as retaliation or unlawful discharge, and does not give them a voice in who owns the employer. Consumer and safety regulation addresses defined harms and has little to say about the governance of speech. Competition law examines the market effects of combinations and was not designed to ask whether the person acquiring a platform ought to hold the capacities it contains. The division of labour assumed that the discarded variables had owners elsewhere, and in this case several of them did not.

Lipton's observation that the market was meant to reinforce the legal protections fits the same pattern. Markets discipline an owner through exit, a point made at length by Hirschman, who distinguished the exit available to those who can leave from the voice available to those who wish to repair what they stay with \cite{hirschman1970}. The sale of Twitter reversed the usual relation between the two. The party that could exit cleanly, the departing shareholders, received full value and had no reason to care about what followed. The parties for whom exit was costly, the users who had accumulated audiences and the employees who depended on their positions, had neither the exit nor a channel of voice that the transaction recognized. The market mechanism that was supposed to price the harm to these parties could not do so, because the only parties whose willingness to pay mattered at the point of sale were those who would not live with the result.

\section{Authorization and Its Remainder}

The distinction between authorization and its remainder clarifies why enforcement of the bargain, however vigorous, cannot repair the omission. To authorize a transition is to certify that it satisfies the conditions the authorizing procedure was built to check. Authorization is necessary, since a transition that fails the check should not proceed, but it does not exhaust the obligations that the transition generates. A transfer of a platform generates obligations of stewardship, of continuity for those who depend on the platform, and of accountability for the decisions the new owner will make. These obligations are the remainder of the authorization. They cannot be discharged by pointing to the fact that the sale was properly approved, and they cannot inherit authority from the approval, because the approval was given by parties who did not have these obligations in view.

Three decisions should be kept distinct here. The board decided to accept the offer, and in doing so was applying the shareholder projection as the law directs. Lipton's account of that direction is exact. Under \emph{Revlon} \cite[182]{revlon1986}, a board operating in the ordinary course may consider employees, customers, and other stakeholders on the expectation that doing so serves the enterprise, but a board pursuing a sale is prohibited from weighing them, because their welfare is no longer plausibly tied to that of the shareholders. The company in this setting is, in her phrase, an avatar for shareholder wealth \cite[393]{lipton2024}, and the merger agreement stated that it had no third-party beneficiaries, so that even the shareholders, who were the substantive beneficiaries, held no enforcement rights of their own. The contractual right to enforce the agreement sat with the company, the substantive interest in the price sat with the shareholders, who also had a recognized role in approving the sale, and a third group, with the largest stake in how the transferred powers would be exercised, held none of these. Three questions are thereby separated: who could approve the sale, who could enforce the agreement, and who could contest the terms on which the transferred powers would be exercised. Shareholders had standing with respect to the terms on which their interests were sold. The contractual right to enforce the agreement belonged to Twitter. Those who would remain subject to the platform's powers had no corresponding standing to contest the terms on which those powers were transferred, and that is the question the procedure did not ask. Twitter then sued to enforce the agreement when the buyer attempted to withdraw, which was an exercise of the company's contractual rights. A court might in turn have decided whether to order specific performance, but the parties closed before it ruled on that remedy, so the question of what equitable discretion could have weighed was never answered. None of these three steps asked who, among the people the platform reached, was entitled to demand continuity, to contest later changes in governance, or to hold the new owner to obligations over time. The board's decision, the company's suit, and the closing each completed a step in the transition without generating a place where that question could be put by someone entitled to put it. The success of the bargain on the retained variables and its silence on the omitted ones were the same fact seen from two sides.

It might be objected that a court could be permitted to weigh harms to outsiders when deciding whether to order a buyer to close. Lipton takes this seriously. She notes that a Delaware court asked for specific performance exercises equitable discretion, that an earlier decision of then-Vice Chancellor Strine attended to the effects of a forced merger on employees, and that there was a real chance the position of nonshareholders would have played a role in the Chancellor's decision. She also treats the reform of remedies as unlikely to change much, for two reasons. If courts weighed externalities more heavily, parties would negotiate large termination fees in advance, and a buyer of Musk's kind would still prefer to complete the purchase, because outside investors would fund part of it and a settlement would leave him with nothing to show for the same capital. And remedies matter only once a deal has been breached, so they do nothing about ill-advised acquisitions that proceed without dispute. The structural reading identifies a further point about what would remain unresolved even if a court had taken up the question. A court that heard evidence about the consequences of the sale for users and employees would have made those consequences visible to the procedure, and it would still have left open whether those groups had a recognized way to represent their interests, to obtain enforceable conditions, or to seek review after closing. A richer projection is not the same as a procedure in which the affected parties are entitled to appear.

\section{Capital Without a Return on the Asset}

The analysis so far concerns the transition. A second question concerns how the buyer came to be in a position to make it, and here the relevant frame is process rather than state. A fortune is the current value of a long history of exchanges, advantages, and enforcement events, and to ask how it arose is to ask which constraints were operating at each stage and how they bore on the one who accumulated it. Lipton argues that weak or delayed enforcement in areas such as workplace safety, vehicle safety, and securities disclosure plausibly helped Musk amass and retain his fortune, and she is careful to say that the counterfactual cannot be known. The caution is appropriate, and the structural point survives it. When enforcement is slow relative to the pace of accumulation, a cost that the rules were meant to impose is deferred, and deferral shifts the incidence of the risk and the cost onto parties who did not agree to bear them.

The financing of the acquisition illustrates a related feature, and it should be described with the care that the evidence permits. Lipton observes that the buyer's wealth gave him the resources to fund much of the purchase himself and also let him raise additional capital from a varied group of lenders and co-investors, many of whom contributed relatively small sums. She takes this to suggest that they did not expect to profit from the Twitter deal so much as to ingratiate themselves with him for the sake of later ones. That is an inference from the size of the contributions to the motives of many of the participants, and the expected return on the asset could still have influenced the financing alongside it. The evidence does not establish how much of any investor's commitment traced to one motive or the other. She adds that the buyer appears to have acquired a public company with no business plan at all, let alone one to increase its profitability, whereas wealth-destroying and socially disruptive acquisitions are ordinarily preceded by at least a plausible plan that the acquirer believes will generate profits, and such plans are assumed to be constrained by market rules.

\begin{proposition}
To the extent that a financier's expected return on a transaction is realized through a relation to a party and not through the asset transacted, the price signal attached to the asset constrains that financier's participation less, and the transaction may proceed at prices that the asset alone would not support.
\end{proposition}

The proposition is conditional on the extent of the relational motive, and the case supplies a reason to think the extent was not zero for many of the smaller participants without supplying a measure of it. It also should not be read as a claim about the shareholder price. The \$54.20 figure was set when the offer was accepted, and Lipton attributes its character as a windfall to the later collapse of the market, so the gain to shareholders cannot be traced to the relational motive of the financiers. The point that survives is narrower. Corporate law counts the successful enforcement of a premium as success, and the premium at the time of enforcement was largely a product of market movement unrelated to any assessment of what the platform would become in the buyer's hands. A procedure that measures its own success by the gap between the agreed price and the standalone value therefore measures something about timing and nothing about stewardship.

\section{A Loop and Not an Event}

The reading offered above treats wealth, enforcement, and control as stages of a possible feedback process, and the case supports some links of that process more firmly than others. The firmly supported links are two. Weak or delayed enforcement can facilitate accumulation, because the costs that enforcement would have imposed are deferred or not borne, and accumulated wealth can purchase control over institutions that shape public communication, as the acquisition itself shows. A third link is possible but is not established by the case. Ownership of a communications platform could in principle weaken oversight of its owner, because the platform controls a channel through which the public learns of conduct it might otherwise demand be regulated. That mechanism belongs in the analysis as a hypothesis about how the loop could close, and it would need evidence of its operation before it could be asserted as a finding.

Seen as a sequence of this kind, the problem cannot be located at any single point, and this explains why the corporate law that governed the sale can be described as functioning flawlessly. Each component works within its own projection. Enforcement agencies act at the pace their resources and legal standards allow. Investors seek returns by their own accounting. Courts enforce valid contracts. The acquisition shows how wealth accumulated under one set of constraints can confer control over public communications under another. If that control also weakens oversight of its holder, the sequence becomes a feedback loop, with no institution in the sequence charged with assessing its cumulative effects. Lipton's title compresses the structural point into a sentence, and the sentence should be read as a claim about composition: each part may satisfy its own constraints while no part is responsible for the constraints that a public would want the whole to meet. Whether the loop closes in this case is a causal question that the case leaves open.

\section{What a Legitimate Transfer Would Require}

The question of legitimacy can now be stated without moralizing about the buyer. A transfer of control over an institution that shapes public life is legitimate to the extent that the procedure authorizing it contains the standing relations appropriate to the consequences. The requirement is not that the public vote on every sale. It is that the procedure be built so that the variables on which the public's stake depends are represented by parties entitled to represent them, with a right to participate, to seek enforceable conditions, and to request review, while the authorizing body holds the power to impose conditions or to refuse approval under stated criteria. Three conditions follow from the preceding analysis.

The grounds on which standing is assigned can be of two kinds, and the argument needs both to be kept apart. An instrumental ground says that a procedure which represents the affected parties will make better decisions, because it will see consequences it would otherwise miss. That ground is the one the informational result supports. Three aims should be kept apart. Visibility satisfies the informational aim of making consequences available for consideration. Participation can serve the aim of better decisions, since parties who are present can supply evidence and expose mistaken assessments, though the size of that benefit is uncertain. Entitlement is a separate ground altogether, which gives a reason for participation even when its informational benefit is uncertain, and neither of the first two supplies it, because a procedure can see a consequence, and can even benefit from the presence of those who bear it, without any party being owed a place in it. A ground of entitlement says that a party is owed a place in the procedure because of its relation to the consequence, whether or not its presence improves the outcome. Three features of the relation support entitlement in this case. The capacities that change hands are exercised over the affected parties and not merely in their vicinity. Their exposure to the new controller begins at closing and has no natural end, whereas the stake of the party that was heard concluded at closing. And no other procedure holds their claim, so that the omission is not covered elsewhere. These features connect to the conditions below because each condition prevents a way in which an entitlement that exists could be defeated: scope matching prevents the population heard from diverging from the population governed, remainder assignment prevents the obligations from falling to no one, and revisability prevents entitlement from lapsing as the consequences unfold. The premise that connects these facts to standing can be stated directly: people subject to a durable power have a claim to contest the terms on which that power is transferred when existing institutions give them no effective avenue for doing so. The premise supports a claim to representation, to enforceable conditions, and to review. It does not support a claim to veto the sale, and the distinction matters because the conditions below are designed to shape the transfer and not to prevent it. This premise is a substantive commitment and it is not derived from the informational result. A reader who rejects it is left with pure instrumentalism, under which visibility, or participation where it improves decisions, suffices and no one is owed anything, and the essay's conclusion then weakens to the claim that the procedure was informationally incomplete. A reader who rejects stakeholder theory in general may still accept the narrower premise, since it depends on the durability of the power and the absence of another avenue and does not assign weight to stakeholders as such.

The conditions can now be stated. The first is scope matching. The population entitled to be heard should correspond to the population over which the transferred capacities are exercised. For a platform of this reach the relevant population includes users, employees, and the organizations that depend on the platform for communication, and the representation of that population need not be direct, so long as it is real. The second is remainder assignment. The obligations that the transition generates beyond the sale itself should be assigned to an identifiable party with the authority to discharge them, and the assignment should be made at the time of authorization and not discovered afterward. A buyer who acquires a capacity should acquire the associated obligations by the same act, in a form that can be checked. The third is revisability. Because the consequences of control over a communications platform unfold over years, the authorization should be subject to review under conditions that were stated in advance, so that the procedure remains sufficient under later revisions of the environment and not only at the moment of closing.

None of these conditions requires a prohibition on private ownership of platforms, and none presupposes that the actual owner will behave badly. Together, these conditions require a procedure that can identify the consequences of the transfer, gives affected parties a recognized means of contesting them, assigns enforceable duties to the holder of the transferred powers, and specifies when compliance can be reviewed. Where wealth is so concentrated that a single party can acquire such an institution without any procedure of this kind, the failure of policy that Lipton identifies is the absence of the procedure, and the wealth is the condition that makes the absence consequential.

\section{Limits of the Argument}

Several limits should be stated. The account of the earlier accumulation rests on Lipton's own qualification that the counterfactual is unknowable, and the structural claim about deferred costs does not depend on any particular attribution of the fortune to any particular failure of enforcement. The treatment of the financing relies on her inference from the small size of many contributions to the motives of those lenders and co-investors, and is held here with the same tentativeness. The conditions proposed in the previous section are conditions on a procedure, and they leave open which institutions could satisfy them, how representation of a diffuse population would be constituted without capture, and what costs a more demanding procedure would impose on transactions that are benign. These questions are genuine, and a proposal that ignored them would reproduce the mistake of the projection it criticizes, by retaining the variables that are convenient and discarding the rest. Finally, the argument is directed at a division of labour and not at the individuals who work within it. Judges applying the doctrine, directors discharging their duties, and shareholders accepting a premium acted within roles that the law defined, and the structural reading explains how conduct that is correct in each role can compose into an outcome that none of the roles was designed to examine.

\section{Conclusion}

The sale of Twitter was a success for corporate law and a failure of policy, and the two descriptions are consistent because they apply to different projections of the same transition. The projection used by the law retained price, premium, and enforceability, and on these it was correct. The projection that the public needed would have retained the capacities that changed hands and the population over which they are exercised, and on these the law was silent. Sufficiency was assumed on the strength of a division of labour that did not hold for a transition of this size. Seeing a consequence and entitling someone to contest it are separate achievements, and the case shows how the first can be approximated while the second is absent. The board could secure a premium, the shareholders could accept it, and Twitter could enforce its contractual rights, and all of this left unresolved who could demand continuity, contest changes in platform governance, or hold the new owner to obligations over time. Authority to complete a sale and legitimacy to exercise the powers it transfers are distinct, and the first does not supply the second. That conclusion holds even if the buyer's motives and the effects of earlier enforcement remain uncertain, since it rests on the structure of the procedure and on neither. What the case demonstrates is that a legal system can be locally sound and globally unanswerable, and that the repair lies in constructing procedures in which the affected parties hold standing, with the obligations of stewardship and continuity assigned at closing, when the capacities that generate them change hands.

\begin{thebibliography}{9}

\bibitem{berle1931}
Berle, A. A., Jr. ``Corporate Powers as Powers in Trust.'' \textit{Harvard Law Review} 44 (1931): 1049.

\bibitem{dodd1932}
Dodd, E. Merrick, Jr. ``For Whom Are Corporate Managers Trustees?'' \textit{Harvard Law Review} 45 (1932): 1145.

\bibitem{dodge1919}
\textit{Dodge v.\ Ford Motor Co.}, 204 Mich.\ 459, 170 N.W.\ 668 (1919).

\bibitem{revlon1986}
\textit{Revlon, Inc.\ v.\ MacAndrews \& Forbes Holdings, Inc.}, 506 A.2d 173 (Del.\ 1986).

\bibitem{friedman1970}
Friedman, Milton. ``The Social Responsibility of Business Is to Increase Its Profits.'' \textit{The New York Times Magazine}, September 13, 1970.

\bibitem{hirschman1970}
Hirschman, Albert O. \textit{Exit, Voice, and Loyalty: Responses to Decline in Firms, Organizations, and States}. Cambridge, MA: Harvard University Press, 1970.

\bibitem{lipton2024}
Lipton, Ann M. ``Every Billionaire Is a Policy Failure.'' \textit{Virginia Law \& Business Review} 18 (2024): 327. Available at SSRN: https://ssrn.com/abstract=4442029.

\end{thebibliography}

\end{document}
